\section{Supplementary derivations}
\label{app:proofs}

\subsection{The cubic form of the closure}

At $q=3/2$, substituting $u=\sqrt m>0$ into \cref{eq:closure} gives
\begin{equation}
  \beta u^3 - \alpha u^2 + \Mzero = 0 .
  \label{eq:cubic}
\end{equation}

\begin{proposition}
\Cref{eq:cubic} has exactly one negative real root for all
$\alpha,\beta,\Mzero>0$. Its two remaining roots are real and positive if and
only if $\Mzero\le\Mmax$, in agreement with \cref{thm:trichotomy}.
\end{proposition}

\begin{proof}
Write $\Pi(u) = \beta u^3-\alpha u^2+\Mzero$. By Vieta, the product of the
three roots is $-\Mzero/\beta<0$, so either one or three roots are negative.
Since $\Pi(0)=\Mzero>0$ and $\Pi(u)\to+\infty$ as $u\to+\infty$, and
$\Pi'(u)=3\beta u^2-2\alpha u = u(3\beta u - 2\alpha)$ has roots $u=0$ and
$u=2\alpha/(3\beta)$, the function decreases on $(0,2\alpha/(3\beta))$ and
increases thereafter, so it has at most two positive roots. Hence exactly one
root is negative. The local minimum value is
\[
  \Pi\!\left(\frac{2\alpha}{3\beta}\right)
  = \beta\frac{8\alpha^3}{27\beta^3} - \alpha\frac{4\alpha^2}{9\beta^2} + \Mzero
  = \frac{8\alpha^3-12\alpha^3}{27\beta^2} + \Mzero
  = \Mzero - \frac{4\alpha^3}{27\beta^2}
  = \Mzero - \Mmax ,
\]
using \cref{eq:closedform}. Two positive real roots exist iff this is
non-positive, i.e.\ iff $\Mzero\le\Mmax$. Note also that
$u^* = 2\alpha/(3\beta)$ satisfies $(u^*)^2 = 4\alpha^2/(9\beta^2) = m^*$, as
required for consistency.
\end{proof}

\Cref{fig:cubic} draws $\Pi$ in the three cases.

\begin{figure}[tbp]
\centering
% Pi(u) = beta u^3 - alpha u^2 + M0 for the worked example, u = sqrt(m).
% u* = 2 alpha / (3 beta) = 1.8688; Pi(u*) = M0 - M0max.
% Roots: M0 = 0.45: -0.7471, 1.0689, 2.4813; M0 = M0max: -0.9344, 1.8688 (double);
% M0 = 0.95: -1.0429 and a complex pair.
\begin{tikzpicture}
\begin{axis}[paperplot, height=6.2cm,
  xlabel={$u=\sqrt m$ [\si{\kilo\gram^{1/2}}]},
  ylabel={$\Pi(u)=\beta u^3-\alpha u^2+\Mzero$},
  xmin=-1.25, xmax=2.95, ymin=-0.52, ymax=1.25, domain=-1.25:2.95,
  samples=200, legend style={at={(0.5,1.03)}, anchor=south}, legend columns=3]
  \fill[cG!10] (axis cs:-1.25,-0.52) rectangle (axis cs:0,1.25);
  \node[lbl, cG] at (axis cs:-0.62,1.12) {$u<0$: not a mass};
  \draw[black, line width=0.5pt] (axis cs:-1.25,0) -- (axis cs:2.95,0);
  \addplot[cC] {0.2271*x^3 - 0.6366*x^2 + 0.45};
  \addlegendentry{$\Mzero=0.45<\Mmax$}
  \addplot[cB] {0.2271*x^3 - 0.6366*x^2 + 0.7411};
  \addlegendentry{$\Mzero=\Mmax$}
  \addplot[cD] {0.2271*x^3 - 0.6366*x^2 + 0.95};
  \addlegendentry{$\Mzero=0.95>\Mmax$}
  \draw[construct] (axis cs:1.8688,-0.52) -- (axis cs:1.8688,1.25);
  \node[lbl, anchor=north west] at (axis cs:1.89,1.2) {$u^*=\tfrac{2\alpha}{3\beta}$};
  \addplot[only marks, mark=*, mark size=1.7pt, cC, forget plot]
    coordinates {(-0.7471,0) (1.0689,0) (2.4813,0)};
  \addplot[only marks, mark=*, mark size=1.9pt, cB, forget plot]
    coordinates {(-0.9344,0) (1.8688,0)};
  \addplot[only marks, mark=*, mark size=1.7pt, cD, forget plot]
    coordinates {(-1.0429,0)};
  \addplot[only marks, mark=square*, mark size=1.6pt, cC, forget plot]
    coordinates {(1.8688,-0.2911)};
  \addplot[only marks, mark=square*, mark size=1.6pt, cD, forget plot]
    coordinates {(1.8688,0.2089)};
  \node[lbl, anchor=north, cC!70!black] at (axis cs:1.8688,-0.32)
    {$\Pi(u^*)=\Mzero-\Mmax<0$};
  \node[lbl, anchor=south, cD] at (axis cs:1.8688,0.24) {$\Pi(u^*)>0$};
  \node[lbl, anchor=south east] at (axis cs:1.06,0.02) {$\sqrt{m_-}$};
  \node[lbl, anchor=south west] at (axis cs:2.49,0.02) {$\sqrt{m_+}$};
  \node[lbl, anchor=north] at (axis cs:-0.75,-0.04) {$u_3<0$};
\end{axis}
\end{tikzpicture}

\caption{The cubic \cref{eq:cubic} in $u=\sqrt m$ for the worked example of
\cref{sec:closure}. $\Pi(0)=\Mzero>0$ is a local maximum and $\Pi$ has its
local minimum at $u^*=2\alpha/(3\beta)$, where $(u^*)^2=m^*$ and
$\Pi(u^*)=\Mzero-\Mmax$. One root is always negative and is not a mass. The
two positive roots $\sqrt{m_\pm}$ exist exactly when the minimum is not above
the axis, and they merge at $u^*$ when $\Mzero=\Mmax$.}
\label{fig:cubic}
\end{figure}

\subsection{Derivative of the closure solution with respect to endurance}

Differentiating $F(m(t_f))=\Mzero(t_f)$ implicitly, with
$\beta = \beta_0 t_f$ and $\Mzero = \mfix + \Paux t_f/\eb$,
\begin{equation}
  F'(m)\,\frac{\dd m}{\dd t_f}
  = \frac{\partial \Mzero}{\partial t_f} + m^{3/2}\frac{\partial\beta}{\partial t_f}
  = \frac{\Paux}{\eb} + \beta_0 m^{3/2} ,
\end{equation}
so
\begin{equation}
  \frac{\dd m}{\dd t_f}
  = \frac{\Paux/\eb + \beta_0 m^{3/2}}{\alpha - \tfrac32\beta\sqrt m} .
  \label{eq:dmdt}
\end{equation}
The denominator is $F'(m)$, which vanishes at the fold; \cref{eq:dmdt} therefore
diverges there, quantifying the remark following \cref{thm:bifurcation}. Both
terms in the numerator are positive. The denominator is positive on the light
branch and negative on the heavy branch (\cref{lem:concave}), so the light root
increases with endurance and the heavy root decreases: as $t_f$ grows the two
branches approach each other and merge at the fold.

\subsection{Rotor speed at maximum thrust}

For a fixed-pitch rotor at fixed collective, $T=k_T\Omega^2$, so the speed
required for the sizing thrust $T_{\max}=\lambda mg/N$ relative to hover
$T_h = mg/N$ is
\begin{equation}
  \frac{\Omega_{\max}}{\Omega_h} = \sqrt{\frac{T_{\max}}{T_h}} = \sqrt{\lambda} .
  \label{eq:omegaratio}
\end{equation}
Hence hover occurs at $1/\sqrt\lambda$ of full rotor speed: at $\lambda=1.8$
this is $74.5\%$, which is the figure quoted in \cref{sec:results} and which
also fixes the voltage headroom the pack must provide.

\subsection{Torque authority of the symmetric quadrotor}

Holding total thrust at $mg$ and the other two moments at zero, the largest
moment about body axis $j$ is the value of the linear program
\begin{equation}
  \max\ \tau_j
  \quad\text{s.t.}\quad
  \mat M\vect u = (mg,\ \tau_j\vect e_j),
  \qquad
  0\le u_i\le\Omega_{\max}^2 ,
  \label{eq:authlp}
\end{equation}
and its counterpart for $-\tau_j$; the authority is the smaller of the two.
The engine solves \cref{eq:authlp} directly for every layout. For the $\times$
quadrotor with alternating spins it has a closed form.

Let $a_i = \mat M_{j+1,i}$. For each of the three moment rows $|a_i|=a$ is the
same for every rotor ($k_TL/\sqrt2$ for roll and pitch, $k_Q$ for yaw), and the
sign vector $s_i=\operatorname{sign}a_i$ is orthogonal to the thrust row and to
the other two moment rows. Write $\vect u = \Omega_h^2\vect 1 + \vect d$. The
thrust constraint is $\sum_i d_i=0$, and $\tau_j = a\sum_i s_i d_i = 2aD$ with
$D=\sum_{s_i=+1}d_i=-\sum_{s_i=-1}d_i$. Each of the two increments is at most
$\Omega_{\max}^2-\Omega_h^2$ and each of the two decrements at most $\Omega_h^2$,
so $D\le2\Delta$ with
\begin{equation}
  \Delta = \min\{\Omega_h^2,\ \Omega_{\max}^2-\Omega_h^2\} .
\end{equation}
The bound is attained by $\vect d=\Delta\vect s$, which satisfies every
constraint by the orthogonality, hence
\begin{equation}
  \tau_j^{\max} = 4a\Delta = \Delta\sum_{i=1}^{4}\big|\mat M_{j+1,i}\big| ,
  \label{eq:authderiv}
\end{equation}
which is \cref{eq:authority}. There is no factor $1/N$: the thrust sum is held by
pairing increments with decrements, not by dividing the moment. For
$\lambda\le2$, $\Delta=(\lambda-1)\Omega_h^2$. \Cref{fig:differential} shows
the two cases of the minimum.

\begin{figure}[tbp]
\centering
% The balanced differential that attains the roll authority of the x
% quadrotor. Rotor i sits at psi_i = pi/4 + i pi/2, i = 0..3 (numbered 1..4),
% so the roll row y_i k_T has signs (+,+,-,-). Squared speeds normalised by
% Omega_h^2; Omega_max^2 = lambda Omega_h^2, Delta = min(1, lambda - 1).
\begin{tikzpicture}[font=\small]
  % --- the quadrotor, top view, body x to the right, y up ------------------
  \begin{scope}[shift={(0,0)}, scale=0.82]
    \draw[line width=0.8pt, cG] (-1.1,-1.1) -- (1.1,1.1);
    \draw[line width=0.8pt, cG] (-1.1,1.1) -- (1.1,-1.1);
    \foreach \ang/\n/\s/\c in {45/1/+/cA, 135/2/+/cA, 225/3/-/cD, 315/4/-/cD}{
      \draw[\c, line width=0.9pt, fill=\c!12] (\ang:1.556) circle (0.52);
      \node[font=\footnotesize] at (\ang:1.556) {$\n$};
      \node[\c, font=\normalsize] at (\ang:2.42) {$\s$};
    }
    \draw[vec] (0,0) -- (0.85,0) node[right, lbl] {$\vect b_1$};
    \draw[vec] (0,0) -- (0,0.85) node[above, lbl] {$\vect b_2$};
  \end{scope}
  \node[lbl, align=center] at (0,2.45)
    {$s_i=\operatorname{sign}\mat M_{2,i}=(+,+,-,-)$\\$\vect s\perp$ thrust, pitch, yaw rows};
  \node[lbl, align=center, cA] at (0,-2.3)
    {$y_i>0$ pair up, $y_i<0$ pair down:\\$\tau_1=\sum_i y_iT_i>0$};
  % --- bars: lambda = 1.8 --------------------------------------------------
  \begin{scope}[shift={(3.0,-1.7)}, yscale=1.35, xscale=0.8]
    \draw[->] (-0.3,0) -- (4.3,0);
    \draw[->] (-0.3,0) -- (-0.3,2.85) node[above, lbl] {$u_i/\Omega_h^2$};
    \foreach \y in {0,1,2}{\draw (-0.35,\y) -- (-0.25,\y) node[left=2pt, lbl] {$\y$};}
    \draw[cD, dashed] (-0.3,1.8) -- (4.3,1.8) node[right, lbl, cD] {$\Omega_{\max}^2$};
    \draw[cG, densely dotted] (-0.3,1) -- (4.3,1) node[right, lbl, cG] {$\Omega_h^2$};
    \foreach \x/\h/\c in {0.5/1.8/cA, 1.5/1.8/cA, 2.5/0.2/cD, 3.5/0.2/cD}{
      \fill[\c!55] (\x-0.3,0) rectangle (\x+0.3,\h);
      \draw[\c] (\x-0.3,0) rectangle (\x+0.3,\h);
    }
    \foreach \x/\n in {0.5/1,1.5/2,2.5/3,3.5/4}{\node[below, lbl] at (\x,0) {$\n$};}
    \draw[dim] (1.95,1) -- (1.95,1.8) node[midway, right, lbl, black] {$\Delta$};
    \draw[dim] (3.95,0.2) -- (3.95,1) node[midway, right, lbl, black] {$\Delta$};
    \node[lbl, align=center] at (2,-0.72)
      {$\lambda=1.8$: $\Delta=\Omega_{\max}^2-\Omega_h^2$,\\increments reach the limit};
  \end{scope}
  % --- bars: lambda = 2.5 --------------------------------------------------
  \begin{scope}[shift={(8.6,-1.7)}, yscale=1.35, xscale=0.8]
    \draw[->] (-0.3,0) -- (4.3,0);
    \draw[->] (-0.3,0) -- (-0.3,2.85) node[above, lbl] {$u_i/\Omega_h^2$};
    \foreach \y in {0,1,2}{\draw (-0.35,\y) -- (-0.25,\y) node[left=2pt, lbl] {$\y$};}
    \draw[cD, dashed] (-0.3,2.5) -- (4.3,2.5) node[right, lbl, cD] {$\Omega_{\max}^2$};
    \draw[cG, densely dotted] (-0.3,1) -- (4.3,1) node[right, lbl, cG] {$\Omega_h^2$};
    \foreach \x/\h/\c in {0.5/2/cA, 1.5/2/cA}{
      \fill[\c!55] (\x-0.3,0) rectangle (\x+0.3,\h);
      \draw[\c] (\x-0.3,0) rectangle (\x+0.3,\h);
    }
    \foreach \x in {2.5,3.5}{\draw[cD, line width=1.4pt] (\x-0.3,0) -- (\x+0.3,0);}
    \foreach \x/\n in {0.5/1,1.5/2,2.5/3,3.5/4}{\node[below, lbl] at (\x,0) {$\n$};}
    \draw[dim] (1.95,1) -- (1.95,2) node[midway, right, lbl, black] {$\Delta$};
    \node[lbl, align=center] at (2,-0.72)
      {$\lambda=2.5$: $\Delta=\Omega_h^2$,\\decrements reach zero};
  \end{scope}
\end{tikzpicture}

\caption{The balanced differential that attains the torque authority of the
$\times$ quadrotor, drawn for roll. The sign vector $\vect s$ of the roll row
is orthogonal to the thrust, pitch and yaw rows, so
$\vect u=\Omega_h^2\vect 1+\Delta\vect s$ changes the roll moment and nothing
else. $\Delta$ is the smaller of the room to speed rotors up and the room to
slow them down: at $\lambda=1.8$ the increments reach $\Omega_{\max}^2$ first,
at $\lambda=2.5$ the decrements reach zero first. Either way each of the four
rotors contributes $\Delta$, which is why \cref{eq:authderiv} carries no factor
$1/N$.}
\label{fig:differential}
\end{figure}

The closed form uses the equal-magnitude and orthogonality properties of this
layout and should not be extended blindly. For the hexacopter with
$\psi_0=\pi/6$ the roll row has $\sum_i|\sin\psi_i| = 4$ and the pitch row
$\sum_i|\cos\psi_i|=2\sqrt3$, because two rotors sit on the pitch axis and
contribute nothing to pitch; the magnitudes are unequal and the LP must be
solved. The same holds for any layout with a failed rotor.

\subsection{Sound power to pressure with the reference constants}

The reference intensity implied by $p_0=\SI{20}{\micro\pascal}$ is
$I_0 = p_0^2/(\rho_0 c_0)$. At $\rho_0=\SI{1.225}{\kilo\gram\per\metre\cubed}$
and $c_0=\SI{343}{\metre\per\second}$,
$\rho_0 c_0 = \SI{420}{\pascal\second\per\metre}$ and
$I_0 = (2\times10^{-5})^2/420 = \SI{9.52e-13}{\watt\per\metre\squared}$, against
the nominal $W_0/\SI{1}{\metre\squared} = \SI{1e-12}{\watt\per\metre\squared}$.
The discrepancy is $10\log_{10}(1/0.952)=\SI{0.21}{\decibel}$, which is why
\cref{eq:lwlp} is stated as an equality: the reference constants are chosen to
make it one to within a fifth of a decibel at standard conditions.
